\section{Podsumowanie i wnioski WIP}
\label{sec:Podsumowanie}

\subsection{Realizacja celów pracy}

Niniejsza praca stawiała sobie za cel zaprojektowanie, implementację i ocenę zautomatyzowanego pipeline CI/CD dla modeli uczenia maszynowego. Poniżej zestawiono cele sformułowane we wprowadzeniu z oceną stopnia ich realizacji.

\begin{itemize}
  \item \textbf{Zaprojektowanie architektury systemu automatyzacji} --- % [UZUPEŁNIĆ: ocena realizacji]
  \item \textbf{Implementacja pipeline CI/CD integrującego walidację danych, trening i wdrożenie} --- % [UZUPEŁNIĆ: ocena realizacji]
  \item \textbf{Wprowadzenie obiektywnej bramki jakości opartej na metrykach ML} --- % [UZUPEŁNIĆ: ocena realizacji]
  \item \textbf{Porównanie podejścia ręcznego z zautomatyzowanym} --- % [UZUPEŁNIĆ: ocena realizacji]
  \item \textbf{Realizacja w ramach budżetu AWS Academy ($\sim$\$50)} --- % [UZUPEŁNIĆ: ostateczny koszt]
\end{itemize}

\subsection{Wnioski}

Na podstawie przeprowadzonych eksperymentów i obserwacji sformułowano następujące wnioski. Kontekst porównawczy oparty jest na literaturze z zakresu DevOps \cite{humble2010continuous} oraz MLOps \cite{kreuzberger2023mlops}.

\begin{enumerate}
  \item Automatyzacja pipeline CI/CD dla modeli ML istotnie skraca czas potrzebny do wdrożenia nowej wersji modelu w porównaniu do podejścia ręcznego. % [UZUPEŁNIĆ: konkretne liczby z eksperymentów]

  \item Wprowadzenie dynamicznej bramki jakości opartej na metrykach MLflow eliminuje ryzyko przypadkowego wdrożenia modelu gorszego od aktualnie produkcyjnego, co jest szczególnie istotne przy niedeterministycznym procesie trenowania.

  \item Wersjonowanie danych za pomocą DVC \cite{dvc_docs} zapewnia pełną odtwarzalność eksperymentów --- dla dowolnej wersji modelu możliwe jest odtworzenie dokładnego zbioru danych, na którym był trenowany.

  \item Zastosowanie k3s zamiast zarządzanego Kubernetes (EKS) okazało się wystarczające dla projektu proof of concept i pozwoliło znacząco ograniczyć koszty infrastruktury.

  \item Architektura stop/start instancji EC2 jest efektywna kosztowo, ale wprowadza opóźnienie na uruchomienie instancji i runnera ($\sim$3--5 minut). W środowisku produkcyjnym z wymaganiami SLA należałoby rozważyć utrzymywanie instancji aktywnej.
\end{enumerate}

\subsection{Ograniczenia}

Projekt realizowany był jako proof of concept z następującymi świadomymi ograniczeniami:

\begin{itemize}
  \item \textbf{Budżet AWS Academy} --- sesje laboratoryjne wygasają co 4 godziny, IAM role mają ograniczone uprawnienia (m.in. brak \texttt{s3:PutObject} dla lokalnych poświadczeń). Wymusiło to przeniesienie wszystkich operacji S3 do GitHub Actions.
  \item \textbf{Single-node k3s} --- brak wysokiej dostępności. Awaria instancji EC2 powoduje niedostępność całego systemu.
  \item \textbf{Mały zbiór danych} --- eksperyment przeprowadzono na zbiorze % [UZUPEŁNIĆ: liczba] obrazów. Zachowanie systemu przy znacznie większych zbiorach nie było testowane.
  \item \textbf{Jeden model} --- architektura jest zaprojektowana pod konkretny model (U-Net dla segmentacji). Generalizacja na inne typy modeli wymagałaby modyfikacji skryptów walidacji i bramki jakości.
\end{itemize}

\subsection{Kierunki dalszego rozwoju}

System stanowi solidną bazę, którą można rozbudować w kilku kierunkach:

\begin{itemize}
  \item \textbf{Multi-node Kubernetes} --- zastąpienie single-node k3s klastrem zapewniającym wysoką dostępność i możliwość skalowania poziomego.
  \item \textbf{Experiment tracking} --- rozbudowa wykorzystania MLflow o śledzenie hiperparametrów i automatyczne porównywanie wielu eksperymentów.
  \item \textbf{Data drift monitoring} --- dodanie mechanizmu wykrywającego rozbieżność między danymi treningowymi a produkcyjnymi.
  \item \textbf{A/B testing} --- wdrożenie mechanizmu stopniowego rollout nowej wersji modelu z automatycznym rollback przy pogorszeniu metryk.
  \item \textbf{Generalizacja} --- abstrakcja pipeline pod inne typy modeli i zadań ML, nie tylko segmentację obrazów.
\end{itemize}
